\section{Conclusão}\label{sec:conclusao}

Foram desenvolvidos e comparados dois controladores laterais, um on-off de três níveis com histerese e o Stanley, integrados ao PI de velocidade da Parte 1. Com referências geométricas, o Stanley obteve erros RMS abaixo de $1\,\mathrm{cm}$, enquanto o on-off apresentou ciclo-limite com erro RMS de $6$--$8\,\mathrm{cm}$ a $1{,}0\,\mathrm{m/s}$ e de $28$--$38\,\mathrm{cm}$ a $2{,}0\,\mathrm{m/s}$.

Com a referência obtida pela câmera (segmentação HSV, projeção no chão e ajuste local), o Stanley completou a volta de $246\,\mathrm{m}$ em todas as velocidades, inclusive em $2{,}5\,\mathrm{m/s}$, em $99{,}7\,\mathrm{s}$ e com erro verdadeiro RMS de $3{,}6\,\mathrm{cm}$; a redução de velocidade nas curvas custou $8\%$ no tempo e reduziu o erro pela metade. O on-off completou a volta apenas até $0{,}75\,\mathrm{m/s}$ ($328\,\mathrm{s}$) e, acima disso, perdeu a linha ou saiu da pista. Nenhum ensaio terminou em capotamento.

No veículo real, o Stanley com os ganhos do simulador seguiu a faixa a $\approx0{,}63\,\mathrm{m/s}$ com erro estimado de $2$--$3\,\mathrm{cm}$, mas revelou um viés de esterçamento a ser corrigido. Conclui-se que o Stanley é claramente superior para percorrer a pista com rapidez e precisão, enquanto o on-off, embora simples, só é adequado em baixa velocidade.
